\section{Worum es geht}

\textbf{[SETZUNG]} Dieses Dokument entwickelt die Dirac-Gleichung in FFGFT vollständig:
die Clifford-Algebra als fundamentale Struktur, den Übergang von der Standard-Dirac-
zur T0-Formulierung, die Massenelimination (Physik ohne Massenparameter formulierbar),
Spin als topologische Eigenschaft, und die modifizierte Dirac-Gleichung aus dem
Zeitfeld-Lagrangian (A142). Das Dokument ergänzt A075 (Feldtheorie, Schrödinger-Limes)
um die vollständige relativistische Struktur.

\section{Die Clifford-Algebra als fundamentale Basis}

\textbf{[B]} Die Standard-Dirac-Gleichung $(i\gamma^\mu\partial_\mu - m)\psi=0$
verwendet $4\times4$ $\gamma$-Matrizen als Darstellung. In FFGFT ist die fundamentale
Struktur die Clifford-Algebra:
\[
  \mathbf e_\mu\mathbf e_\nu + \mathbf e_\nu\mathbf e_\mu = 2g_{\mu\nu}.
\]
Die $\gamma$-Matrizen sind eine Darstellung -- sie werden nicht beseitigt, sondern
auf ihren geometrischen Ursprung zurückgeführt. Das ist der Unterschied: im SM
sind $\gamma$-Matrizen das Fundament; in FFGFT ist es die Clifford-Algebra, die
$\gamma$-Matrizen sind nur eine Darstellung.

\textbf{[B]} \textbf{Was ändert sich in T0?} Die Standard-Dirac-Gleichung hat
$m$ als festen Parameter. In T0 ist die Masse $m(x,t)$ ein dynamisches Feld:
\[
  \bigl[i\gamma^\mu(\partial_\mu + \Gamma_\mu^{(T)}) - m(x,t)\bigr]\psi = 0,
\]
mit Zeitfeld-Verbindung $\Gamma_\mu^{(T)}=-\partial_\mu m/m$ (A142). Im Grenzfall
$\partial_\mu m\to0$ (konstante Masse) reduziert sich das auf die Standard-Dirac-Gleichung.
}

\section{Spin als topologische Eigenschaft}

\textbf{[B]} In FFGFT ist Spin keine intrinsische Eigenschaft, sondern eine
topologische: Spin-$\tfrac12$ entspricht der Toruswicklung $(n_\theta,n_\phi)=(1,1)$.
Die 720°-Symmetrie (zwei Umdrehungen bringen ein Spinor in den Ausgangszustand) folgt
direkt aus der Topologie des Torus -- sie muss nicht postuliert werden.

\textbf{[B]} Ein relativistischer Spinor-Zustand in der Torusdarstellung:
\[
  \Phi(x,t) = \rho(x,t)\,e^{i\theta(x,t)}\cdot\chi(\sigma),
\]
mit Spin-Komponente $\chi(\sigma)$ aus den geometrischen Knotenmoden (A075).
Die Quantenzahl $j=l\pm\tfrac12$ für Fermionen folgt aus der Wicklungsstruktur.

\section{Massenelimination: Physik ohne Massenparameter}

\textbf{[K]} Im T0-Rahmen können alle Observablen als Verhältnisse formuliert werden,
die keinen absoluten Massenparameter benötigen. Die Grundidee: statt $m$ in SI-Einheiten
als Parameter zu führen, wird alles durch $\xi$ und Zeitintervalle ausgedrückt.

\textbf{[K]} \textbf{Systematische Massenelimination.} Aus der T0-Dualität $\tilde T\cdot m=1$
folgt $m=1/\tilde T$. Einsetzen in die Dirac-Gleichung:
\[
  \bigl[i\gamma^\mu\partial_\mu - 1/\tilde T(x,t)\bigr]\psi = 0.
\]
Der Massenparameter verschwindet aus der Gleichung; was bleibt, ist das Zeitfeld
$\tilde T(x,t)$ als dynamische Größe. Das ist die vollständig massenfreie T0-Formulierung.

\textbf{[S]} \textbf{Philosophische Konsequenz.} Wenn alle Physik in Verhältnissen
formulierbar ist (ohne absolute Massen), dann sind Massen keine fundamentalen
Objekte, sondern Hilfskonstrukte der Messung. Was fundamental ist: Zeitintervalle
und die Geometrie des Torus. Das ist dieselbe These wie A085 (natürliche Einheiten)
in relativistischer Sprache.

\section{Massenproportionale Kopplung}

\textbf{[K]} Im T0-Lagrangian koppelt das Zeitfeld massenproportional:
$g_T^\ell = m_\ell\cdot\xi$. Das ist der Mechanismus, der hinter der g-2-Berechnung
(A138) steht: die Kopplung ist nicht eine freie Yukawa-Konstante, sondern durch die
Masse und $\xi$ vollständig bestimmt. Die Ein-Schleifen-Formel für das anomale
magnetische Moment folgt direkt aus dieser massenproportionalen Kopplung.

\section{Verhältnisse vs. Absolutwerte in der Dirac-Struktur}

\textbf{[B]} Dieselbe Unterscheidung wie in A261 und A138: Verhältnisse von
Dirac-Korrekturen (z.B.\ $a_\mu/a_e$) sind korrekturfrei und folgen aus der Struktur
der Clifford-Algebra und den Quantenzahlen. Absolutwerte der $g-2$-Terme brauchen die
SI-Brückenkonstanten und $K_\text{frak}$ (A138).

\section{Prüfskript}

\begin{itemize}[leftmargin=2em,itemsep=2pt,topsep=2pt]
  \item Numerische Verifikation: Standard-Dirac-Gleichung als Grenzfall der T0-Form
    ($\partial_\mu m\to0$); Spin-$\tfrac12$ aus Wicklung $(1,1)$; massenfreie
    Formulierung für einfache Fälle:
    \texttt{python/A\_Serie\_Skripte/a263\_dirac.py}
\end{itemize}

\section{Was offen bleibt}

\textbf{Die vollständige massenfreie Formulierung für alle Wechselwirkungen fehlt.}
Die Massenelimination ist für freie Felder gezeigt; für Wechselwirkungsterme ist sie
nicht systematisch ausgeführt.

\textbf{Der Quark-Sektor fehlt.} Die Dirac-Gleichung für Quarks mit dynamischer Masse
und Farbladung ist nicht auf demselben Niveau entwickelt wie für Leptonen (A150).

\textbf{Das Verhältnis zur QED-Renormierung ist nicht ausgeführt.} Wie die T0-Dirac-
Gleichung sich zur perturbativen QED verhält (Feynman-Regeln, Schleifenintegrals) ist
nicht systematisch entwickelt (A075 offene Hopf-Algebra).

\section{Ersetzt}

Dieses Dokument nimmt auf: Dok.~051 (T0-Dirac-Gleichung, Clifford-Algebra, Spin als
Topologie, massenproportionale Kopplung, Verhältnisse vs.\ Absolutwerte), Dok.~052
(systematische Massenelimination, vollständige massenfreie T0-Formulierung), Dok.~053
(Massenelimination im Dirac-Lagrangian, Verhältnis-basierte Physik), Dok.~050
(vereinfachte Dirac-Gleichung).

\section{Verwendung von KI-Werkzeugen}

Dieses Dokument ist unter Verwendung KI-gestützter Werkzeuge entstanden. Verantwortung
für Inhalt und Fehler liegt beim Autor. Wortlaut der Regeln in A010.
